% Tuba Nazlı İpekyüz — 2200004532
% Options for packages loaded elsewhere
\PassOptionsToPackage{unicode}{hyperref}
\PassOptionsToPackage{hyphens}{url}
\documentclass[
  10pt,
  a4paper,
]{article}
\usepackage{xcolor}
\usepackage[margin=22mm]{geometry}
\usepackage{amsmath,amssymb}
\setcounter{secnumdepth}{5}
\usepackage{iftex}
\ifPDFTeX
  \usepackage[T1]{fontenc}
  \usepackage[utf8]{inputenc}
  \usepackage{textcomp} % provide euro and other symbols
\else % if luatex or xetex
  \usepackage{unicode-math} % this also loads fontspec
  \defaultfontfeatures{Scale=MatchLowercase}
  \defaultfontfeatures[\rmfamily]{Ligatures=TeX,Scale=1}
\fi
\usepackage{lmodern}
\ifPDFTeX\else
  % xetex/luatex font selection
\fi
% Use upquote if available, for straight quotes in verbatim environments
\IfFileExists{upquote.sty}{\usepackage{upquote}}{}
\IfFileExists{microtype.sty}{% use microtype if available
  \usepackage[]{microtype}
  \UseMicrotypeSet[protrusion]{basicmath} % disable protrusion for tt fonts
}{}
\makeatletter
\@ifundefined{KOMAClassName}{% if non-KOMA class
  \IfFileExists{parskip.sty}{%
    \usepackage{parskip}
  }{% else
    \setlength{\parindent}{0pt}
    \setlength{\parskip}{6pt plus 2pt minus 1pt}}
}{% if KOMA class
  \KOMAoptions{parskip=half}}
\makeatother
\usepackage{listings}
\newcommand{\passthrough}[1]{#1}
\lstset{defaultdialect=[5.3]Lua}
\lstset{defaultdialect=[x86masm]Assembler}
\usepackage{longtable,booktabs,array}
\usepackage{caption}
\captionsetup[table]{skip=6pt}
\newcounter{none} % for unnumbered tables
\usepackage{calc} % for calculating minipage widths
% Correct order of tables after \paragraph or \subparagraph
\usepackage{etoolbox}
\makeatletter
\patchcmd\longtable{\par}{\if@noskipsec\mbox{}\fi\par}{}{}
\makeatother
\setlength{\emergencystretch}{3em} % prevent overfull lines
\providecommand{\tightlist}{%
  \setlength{\itemsep}{0pt}\setlength{\parskip}{0pt}}
% Tuba Nazlı İpekyüz — 2200004532
\usepackage{xcolor}
\usepackage{etoolbox}
\usepackage{needspace}
\usepackage{xurl}
\usepackage{newunicodechar}
\definecolor{codebg}{HTML}{F4F6F8}
\definecolor{accent}{HTML}{174A67}
\AtBeginDocument{\hypersetup{linkcolor=accent,urlcolor=accent}}
\setlength{\emergencystretch}{4em}
\newunicodechar{🧠}{\texttt{[U+1F9E0]}}
\lstdefinelanguage{JavaScript}{morekeywords={const,let,function,return,if,else,for,of,in,try,catch,finally,throw,new,import,from,export,default,class,extends,async,await,yield},sensitive=true,morecomment=[l]{//},morecomment=[s]{/*}{*/},morestring=[b]',morestring=[b]",morestring=[b]`}
\lstdefinelanguage{TypeScript}[]{JavaScript}{morekeywords={type,infer,readonly,unknown,never,as,keyof,interface}}
\lstdefinelanguage{jsx}[]{JavaScript}{}
\lstdefinelanguage{tsx}[]{TypeScript}{}
\lstdefinelanguage{js}[]{JavaScript}{}
\lstdefinelanguage{ts}[]{TypeScript}{}
\lstdefinelanguage{text}{}
\lstdefinelanguage{sh}{}
\lstset{basicstyle=\ttfamily\footnotesize,backgroundcolor=\color{codebg},breaklines=true,breakatwhitespace=false,columns=fullflexible,keepspaces=true,showstringspaces=false,frame=single,framerule=0pt,xleftmargin=4pt,xrightmargin=4pt,aboveskip=8pt,belowskip=8pt,tabsize=2,escapeinside={(*@}{@*)}}
\AtBeginEnvironment{longtable}{\footnotesize}
\setlength{\parskip}{5pt}
\setlength{\parindent}{0pt}



\pretocmd{\section}{\Needspace{14\baselineskip}}{}{}
\pretocmd{\subsection}{\Needspace{14\baselineskip}}{}{}


\AtBeginDocument{\raggedright}


\usepackage{bookmark}
\IfFileExists{xurl.sty}{\usepackage{xurl}}{} % add URL line breaks if available
\urlstyle{same}
% fallback for those not using the hyperref driver hyperxmp:
\makeatletter
\@ifundefined{xmpquote}{\newcommand{\xmpquote}[1]{#1}}{}
\makeatother
\hypersetup{
  pdftitle={Destructuring Field Guide},
  hidelinks,
  pdfcreator={LaTeX via pandoc}}

\title{Destructuring Field Guide}
\author{Tuba Nazlı İpekyüz - 2200004532}
\date{}

\begin{document}
\maketitle

{
\setcounter{tocdepth}{3}
\tableofcontents
}
Name: Tuba Nazlı İpekyüz\\
Student number: 2200004532

Bonuses claimed: \textbf{Bonus 1: LaTeX}, \textbf{Bonus 2: Depth},
\textbf{Bonus 3: Real code}. \textbf{Bonus 4: Git repository} is
included with a public repository, as requested by the instructor.

This guide uses a study-session planner. I used AI help to prepare the
code and explanations. I still need to review and explain the work
myself before submission. The project includes forms 1--12, five
pitfalls, T1--T5 and P1--P7. They were checked.

\section{Run}\label{run}

Requires Node.js 20.19 or newer.

\begin{minipage}{\linewidth}

\begin{lstlisting}[language=sh]
npm install
npm run all
npm run check
npm run diagnose # separate intentional compiler failures
\end{lstlisting}

\end{minipage}\par

\texttt{output.\allowbreak{}txt} saves the real output of this project.
JS and TS print the same values for forms 1--12. \texttt{npm run all}
also runs the review and extra examples.

\clearpage

\section{Field guide}\label{field-guide}

\subsection{Form 1: values by position}\label{form-1-values-by-position}

\textbf{Syntax:} \texttt{const [work, pause] = sessionLengths(40)}

\textbf{Example} (\texttt{js/\allowbreak{}arrays.\allowbreak{}mjs},
typed version in \texttt{ts/\allowbreak{}arrays.\allowbreak{}ts}):

\begin{minipage}{\linewidth}

\begin{lstlisting}
function sessionLengths(workMinutes) {
  return [workMinutes, Math.round(workMinutes / 5)]
}
const [work, pause] = sessionLengths(40)
console.log('form 1:', work, pause)
\end{lstlisting}

\end{minipage}\par

\textbf{Output:} \texttt{form 1:\allowbreak{} 40 8}

\textbf{Use it when:} a function returns a few values in a fixed order.

\textbf{Watch out:} The first name gets the first value. The second name
gets the second value. Names do not choose the positions. Without
destructuring, save the result once and read \texttt{result[0]} and
\texttt{result[1]}. In TS, \texttt{[number, number]} says both items
exist.

\subsection{Form 2: skip and rest}\label{form-2-skip-and-rest}

\textbf{Syntax:}
\texttt{const [, nextTopic, .\allowbreak{}.\allowbreak{}.\allowbreak{}laterTopics] = topics}

\textbf{Example} (\texttt{js/\allowbreak{}arrays.\allowbreak{}mjs}):

\begin{minipage}{\linewidth}

\begin{lstlisting}
const topics = ['arrays', 'objects', 'parameters', 'props']
const [, nextTopic, ...laterTopics] = topics
console.log('form 2:', nextTopic, JSON.stringify(laterTopics))
\end{lstlisting}

\end{minipage}\par

\textbf{Output:}
\texttt{form 2:\allowbreak{} objects ["parameters","props"]}

\textbf{Use it when:} you need one item and all the items after it.

\textbf{Watch out:} Rest must come last. It makes a new array and does
not change topics. For this example, use \texttt{topics[1]} and
\texttt{topics.\allowbreak{}slice(2)} without destructuring. With
\texttt{noUncheckedIndexedAccess}, nextTopic has type
\texttt{string | undefined}. The array may not have a second item.

\subsection{Form 3: swap}\label{form-3-swap}

\textbf{Syntax:} \texttt{[morning, evening] = [evening, morning]}

\textbf{Example} (\texttt{js/\allowbreak{}arrays.\allowbreak{}mjs}):

\begin{minipage}{\linewidth}

\begin{lstlisting}
let morning = 'TypeScript'
let evening = 'JavaScript'
;[morning, evening] = [evening, morning]
console.log('form 3:', morning, evening)
\end{lstlisting}

\end{minipage}\par

\textbf{Output:} \texttt{form 3:\allowbreak{} JavaScript TypeScript}

\textbf{Use it when:} you want to swap two variables.

\textbf{Watch out:} Use let variables because their values must change.
The first semicolon keeps this line separate from the previous line.
Without destructuring, save morning in a temporary variable before
changing it. In TS, the new value must fit the type of the variable that
receives it.

\subsection{Form 4: array defaults}\label{form-4-array-defaults}

\textbf{Syntax:}
\texttt{const [duration = 30, reminder = 'enabled', rounds = 2] = settings}

\textbf{Example} (\texttt{js/\allowbreak{}arrays.\allowbreak{}mjs}):

\begin{minipage}{\linewidth}

\begin{lstlisting}
const settings = [undefined, null]
const [duration = 30, reminder = 'enabled', rounds = 2] = settings
console.log('form 4:', duration, reminder, rounds)
\end{lstlisting}

\end{minipage}\par

\textbf{Output:} \texttt{form 4:\allowbreak{} 30 null 2}

\textbf{Use it when:} a missing array item should use a default value.

\textbf{Watch out:} A default is used for undefined or when no item was
given. It is not used for null. Without destructuring, check each item
with \texttt{=== undefined}. The TS example uses a two-item tuple for
the first two values. It uses a separate empty \texttt{number[]} for
rounds, so it does not read past the end of that tuple.

\subsection{Form 5: properties by name}\label{form-5-properties-by-name}

\textbf{Syntax:} \texttt{let \{ title, status \} = session}

\textbf{Example} (\texttt{js/\allowbreak{}objects.\allowbreak{}mjs},
typed version in \texttt{ts/\allowbreak{}objects.\allowbreak{}ts}):

\begin{minipage}{\linewidth}

\begin{lstlisting}
const session = { id: 's1', title: 'Review closures', status: 'planned' }
let { title, status } = session
console.log('form 5:', title, status)
\end{lstlisting}

\end{minipage}\par

\textbf{Output:} \texttt{form 5:\allowbreak{} Review closures planned}

\textbf{Use it when:} you need several properties from one object.

\textbf{Watch out:} Object patterns use property names. The order of the
properties does not matter. Without destructuring, write
\texttt{let title = session.\allowbreak{}title} and
\texttt{let status = session.\allowbreak{}status}. Changing these string
variables does not change session.

\subsection{Form 6: rename}\label{form-6-rename}

\textbf{Syntax:}
\texttt{const \{ title:\allowbreak{} heading \} = session}

\textbf{Example} (\texttt{js/\allowbreak{}objects.\allowbreak{}mjs},
using the session from form 5):

\begin{minipage}{\linewidth}

\begin{lstlisting}
const { title: heading } = session
console.log('form 6:', heading)
\end{lstlisting}

\end{minipage}\par

\textbf{Output:} \texttt{form 6:\allowbreak{} Review closures}

\textbf{Use it when:} you want a clearer name or the property name is
already in use.

\textbf{Watch out:} This creates heading. It does not create another
variable named title. It means
\texttt{const heading = session.\allowbreak{}title}. The colon changes
the local name in both JS and TS. To add a type in TS, write it after
the whole pattern.

\subsection{Form 7: object defaults}\label{form-7-object-defaults}

\textbf{Syntax:} \texttt{const \{ minutes = 40 \} = options}

\textbf{Example} (\texttt{js/\allowbreak{}objects.\allowbreak{}mjs}):

\begin{minipage}{\linewidth}

\begin{lstlisting}
const { minutes = 40 } = {}
const { minutes: explicitNull = 40 } = { minutes: null }
console.log('form 7:', minutes, explicitNull)
\end{lstlisting}

\end{minipage}\par

\textbf{Output:} \texttt{form 7:\allowbreak{} 40 null}

\textbf{Use it when:} an optional property should use a default value.

\textbf{Watch out:} Defaults replace undefined or missing values, never
null. The source cannot be null. In TS, null stays allowed after a
default. Without destructuring, check \texttt{=== undefined} before
using 40.

\subsection{Form 8: object rest}\label{form-8-object-rest}

\textbf{Syntax:}
\texttt{const \{ id, .\allowbreak{}.\allowbreak{}.\allowbreak{}details \} = session}

\textbf{Example} (\texttt{js/\allowbreak{}objects.\allowbreak{}mjs},
using the session from form 5):

\begin{minipage}{\linewidth}

\begin{lstlisting}
const { id, ...details } = session
console.log('form 8:', id, JSON.stringify(details))
\end{lstlisting}

\end{minipage}\par

\textbf{Output:}
\texttt{form 8:\allowbreak{} s1 \{"title":\allowbreak{}"Review closures","status":\allowbreak{}"planned"\}}

\textbf{Use it when:} you want to keep id separate from the other
properties.

\textbf{Watch out:} Object rest makes a new object. It copies the
remaining own enumerable properties: the object's own properties that
are marked for enumeration. It does not make separate copies of objects
inside those properties. For our plain object, read id separately and
put title and status in a new object. Rest collects properties into an
object. Spread uses those properties in an expression.

\subsection{Form 9: nested patterns}\label{form-9-nested-patterns}

\textbf{Syntax:}
\texttt{const \{ room:\allowbreak{} \{ label \}, topics:\allowbreak{} [firstTopic] \} = plan}

\textbf{Example} (\texttt{js/\allowbreak{}objects.\allowbreak{}mjs}):

\begin{minipage}{\linewidth}

\begin{lstlisting}
const plan = { room: { label: 'Library B' }, topics: ['closures', 'props'] }
const { room: { label }, topics: [firstTopic] } = plan
console.log('form 9:', label, firstTopic)
\end{lstlisting}

\end{minipage}\par

\textbf{Output:} \texttt{form 9:\allowbreak{} Library B closures}

\textbf{Use it when:} you need values inside an object or array.

\textbf{Watch out:} This creates label and firstTopic, not room or
topics. If room is missing, reading label throws an error. Without
destructuring, read \texttt{plan.\allowbreak{}room.\allowbreak{}label}
and \texttt{plan.\allowbreak{}topics[0]}. In TS, firstTopic may still be
undefined because topics can be empty.

\subsection{Form 10: computed property}\label{form-10-computed-property}

\textbf{Syntax:}
\texttt{const \{ [key]:\allowbreak{} selectedValue \} = session}

\textbf{Example} (\texttt{js/\allowbreak{}objects.\allowbreak{}mjs},
using the session from form 5):

\begin{minipage}{\linewidth}

\begin{lstlisting}
const key = 'status'
const { [key]: selectedValue } = session
console.log('form 10:', key, selectedValue)
\end{lstlisting}

\end{minipage}\par

\textbf{Output:} \texttt{form 10:\allowbreak{} status planned}

\textbf{Use it when:} a variable holds the name of the property you
need.

\textbf{Watch out:} \texttt{[key]} uses the value stored in key. Without
brackets, key would mean a property named key. The plain version is
\texttt{const selectedValue = session[key]}. In TS,
\texttt{keyof Session} limits the variable to known property names.

\subsection{Form 11: parameters}\label{form-11-parameters}

\textbf{Syntax:}
\texttt{function makeSession(\{ title = 'Free practice', minutes = 40 \} = \{\})}

\textbf{Example} (\texttt{js/\allowbreak{}params.\allowbreak{}mjs},
typed version in \texttt{ts/\allowbreak{}params.\allowbreak{}ts}):

\begin{minipage}{\linewidth}

\begin{lstlisting}
function makeSession({ title = 'Free practice', minutes = 40 } = {}) {
  return { title, minutes }
}
console.log('form 11:', JSON.stringify(makeSession()))
console.log('form 11 custom:', JSON.stringify(makeSession({ title: 'React', minutes: 25 })))
function totalMinutes([work, pause]) { return work + pause }
console.log('form 11 array:', totalMinutes([40, 8]))
\end{lstlisting}

\end{minipage}\par

\Needspace{7\baselineskip}

\textbf{Output:}

\begin{minipage}{\linewidth}

\begin{lstlisting}
form 11: {"title":"Free practice","minutes":40}
form 11 custom: {"title":"React","minutes":25}
form 11 array: 48
\end{lstlisting}

\end{minipage}\par

\textbf{Use it when:} a function takes named options or a small tuple.

\textbf{Watch out:} The outer \texttt{= \{\}} handles a missing
argument. The inner defaults handle missing properties. Neither replaces
null. Without the pattern, receive \texttt{options = \{\}} and read the
properties inside the function. In TS, put
\texttt{:\allowbreak{} Options} after the pattern. Its properties must
be optional for the default \texttt{\{\}} to work.

\subsection{Form 12: loop entries}\label{form-12-loop-entries}

\textbf{Syntax:}
\texttt{for (const [topic, minutes] of Object.\allowbreak{}entries(targets))}

\textbf{Example} (\texttt{js/\allowbreak{}params.\allowbreak{}mjs}):

\begin{minipage}{\linewidth}

\begin{lstlisting}
const targets = { reading: 20, coding: 40 }
for (const [topic, minutes] of Object.entries(targets)) {
  console.log('form 12 object:', topic, minutes)
}
const completed = new Map([['Monday', 2], ['Tuesday', 3]])
for (const [day, count] of completed) {
  console.log('form 12 map:', day, count)
}
\end{lstlisting}

\end{minipage}\par

\Needspace{7\baselineskip}

\textbf{Output:}

\begin{minipage}{\linewidth}

\begin{lstlisting}
form 12 object: reading 20
form 12 object: coding 40
form 12 map: Monday 2
form 12 map: Tuesday 3
\end{lstlisting}

\end{minipage}\par

\textbf{Use it when:} you need both the key and value in a loop.

\textbf{Watch out:} Use for\ldots of to read entries. A for\ldots in
loop gives property names instead. Without destructuring, store each
pair in entry and read entry{[}0{]} and entry{[}1{]}. The TS example
uses \texttt{Record<string, number>} for the object and
\texttt{Map<string, number>} for the Map.

\section{Pitfalls}\label{pitfalls}

The examples are in \texttt{js/\allowbreak{}pitfalls.\allowbreak{}mjs}.
The invalid syntax is tested with \texttt{new Function} and its error is
caught. This lets the rest of the file run.

\begin{enumerate}
\def\labelenumi{\arabic{enumi}.}
\tightlist
\item
  You cannot take properties from undefined. The actual output was
  \texttt{pitfall 1:\allowbreak{} TypeError:\allowbreak{} Cannot destructure property 'minutes' of 'undefined' as it is undefined.\allowbreak{}}
  Check the source first. A parameter default can handle undefined, but
  not null. You also cannot take properties from null.
\item
  A default does not replace null.
  \texttt{\{ minutes = 40 \} = \{ minutes:\allowbreak{} null \}} printed
  \texttt{pitfall 2:\allowbreak{} null}. Defaults are only used for
  undefined or missing values.
\item
  At the start of a statement, \texttt{\{} is read as a code block. The
  attempt without parentheses printed
  \texttt{pitfall 3 invalid:\allowbreak{} SyntaxError:\allowbreak{} Unexpected token '='}.
  Adding parentheses made it an expression. The fixed version
  \texttt{(\{ minutes:\allowbreak{} assignedMinutes \} = \{ minutes:\allowbreak{} 25 \})}
  printed \texttt{pitfall 3 fixed:\allowbreak{} 25}.
\item
  Two variables can point to the same object. Changing room.label
  changed the original room. Changing the string id in the rest object
  did not change the original id. The result was
  \texttt{pitfall 4:\allowbreak{} Library C s1 s2}. Rest is a shallow
  copy: it does not make a separate copy of the nested room object.
\item
  Object patterns use names. Array patterns use positions. Changing the
  order of object properties still gave
  \texttt{pitfall 5 object:\allowbreak{} React practice planned}. The
  array gave
  \texttt{pitfall 5 array:\allowbreak{} planned React practice} because
  its positions were different.
\end{enumerate}

\section{TypeScript topics T1--T5}\label{typescript-topics-t1t5}

Code: \texttt{ts/\allowbreak{}types.\allowbreak{}ts}. TypeScript checks
types before the code runs. It does not change how JavaScript takes
values from a pattern.

\subsection{T1: a colon can rename}\label{t1-a-colon-can-rename}

\texttt{const \{ minutes:\allowbreak{} string \} = \{ minutes:\allowbreak{} 25 \}}
creates a variable named string. Its value is 25, so its type is number.
The word string here is a name, not a type. To add a type, write it
after the whole pattern:
\texttt{const \{ minutes:\allowbreak{} duration \}:\allowbreak{} \{ minutes:\allowbreak{} number \} = .\allowbreak{}.\allowbreak{}.\allowbreak{}}.
Actual output: \texttt{T1:\allowbreak{} 25 number 25}. Calling
\texttt{string.\allowbreak{}toUpperCase()} gives TS2339 because numbers
do not have that method.

\subsection{T2: tuples and as const}\label{t2-tuples-and-as-const}

\texttt{sessionSummary():\allowbreak{} [number, string]} returns two
items: a number, then a string. This type is called a tuple. It tells
TypeScript the type at each position.
\texttt{const fixed = [25, 'props'] as const} gives the type
\texttt{readonly [25, 'props']}. It keeps those exact values as types
and prevents changing the tuple's positions through this variable. It
does not freeze the array at runtime. Actual output:
\texttt{T2:\allowbreak{} 40 closures 25 props}.

\subsection{T3: unchecked array access}\label{t3-unchecked-array-access}

A \texttt{string[]} can be empty. With noUncheckedIndexedAccess enabled,
\texttt{const [first] = topics} gives first the type
\texttt{string | undefined}. Before calling a string method, check
\texttt{first !== undefined}. Another option is a default:
\texttt{[safeFirst = 'free practice']}. Actual output:
\texttt{T3:\allowbreak{} undefined free practice}. Calling the method
without a check gives TS18048:
\texttt{'first' is possibly 'undefined'.\allowbreak{}}

\subsection{T4: optional fields and
defaults}\label{t4-optional-fields-and-defaults}

The question mark in \texttt{minutes?:\allowbreak{} number} means the
property can be missing. Without a default, the extracted value has type
\texttt{number | undefined}. With \texttt{minutes = 40}, it has type
number. A default is used when the value is undefined or has not been
supplied. It is not used for null. If null is allowed, it stays in the
type and the value. Actual output:
\texttt{T4:\allowbreak{} undefined 40 null}. Calling a number method on
a possibly null value gives TS18047:
\texttt{'minutes' is possibly 'null'.\allowbreak{}}

\subsection{T5: wrong property and fix}\label{t5-wrong-property-and-fix}

The source has minutes, but
\texttt{\{ duration:\allowbreak{} wrongName \}} asks for duration.
TypeScript reports TS2339:
\texttt{Property 'duration' does not exist on type '\{ minutes:\allowbreak{} number; \}'.\allowbreak{}}
The one allowed \texttt{@ts-expect-error} marks this example. The fix is
to ask for the existing property:
\texttt{\{ minutes:\allowbreak{} correctedName \}}. Actual output:
\texttt{T5:\allowbreak{} undefined 25}. The comment lets us demonstrate
the error; it does not fix the property lookup. TypeScript also reports
an error if an expect-error comment is no longer needed.

\section{Props and children P1--P7}\label{props-and-children-p1p7}

Run the examples in \texttt{js/\allowbreak{}props.\allowbreak{}jsx} and
\texttt{ts/\allowbreak{}props.\allowbreak{}tsx}. They use
\texttt{renderToStaticMarkup} to produce HTML as text. The output below
comes from real runs. These runs do not test browser clicks or events.

\subsection{P1: one props object, two ways to read
it}\label{p1-one-props-object-two-ways-to-read-it}

React passes one props object to each component. \texttt{CardA(props)}
reads \texttt{props.\allowbreak{}title} and
\texttt{props.\allowbreak{}children} in the function body.
\texttt{Card(\{ title, children \})} reads the same properties in the
parameter list. This uses form 5 for object properties and form 11 for
function parameters. Both versions receive one object and produce the
same HTML.

\begin{minipage}{\linewidth}

\begin{lstlisting}
function Card({ title, children }) {
  return <section><h2>{title}</h2>{children}</section>
}
\end{lstlisting}

\end{minipage}\par

\subsection{P2: children is a field in that
object}\label{p2-children-is-a-field-in-that-object}

\texttt{<Card title="Study">Practice</\allowbreak{}Card>} and
\texttt{<Card title="Study" children="Practice" /\allowbreak{}>} produce
the same HTML. Children is one property inside props. React does not
pass it as a second argument.

In the JS example,
\texttt{<Card children="Attribute">Nested</\allowbreak{}Card>} displays
Nested. Content between the tags wins when children is also written as
an attribute. TypeScript reports TS2710 for this example. We check that
error separately in \texttt{compiler-diagnostics.\allowbreak{}txt}, so
the normal TS file still passes its check.

\subsection{P3: the value of children
matters}\label{p3-the-value-of-children-matters}

A string appears as text. A number, including 0, also appears as text. A
React element becomes HTML. Several elements appear in order. Undefined,
null and false show no child content. The section and heading still
appear.

React does not call a function child for us. The JS example passes a
function without calling it. React logs a warning and shows no child
content. RenderCard works differently: it calls \texttt{children(25)}
and displays the returned element.

\begin{minipage}{\linewidth}

\begin{lstlisting}
function RenderCard({ children }) {
  return <section>{children(25)}</section>
}
\end{lstlisting}

\end{minipage}\par

For this function child, the TS type is
\texttt{\{ children:\allowbreak{} (minutes:\allowbreak{} number) => ReactNode \}}.
It is different from ordinary
\texttt{children?:\allowbreak{} ReactNode}. The JS file copies React's
warning to stdout so it appears in output.txt. It restores console.error
in finally. The \texttt{\%s} text is part of React's raw warning.

\subsection{P4: separate component props, forward DOM
props}\label{p4-separate-component-props-forward-dom-props}

\begin{minipage}{\linewidth}

\begin{lstlisting}
function RestFirst({ title, children, ...rest }) {
  return <section {...rest} title={title}>{children}</section>
}
\end{lstlisting}

\end{minipage}\par

The pattern takes title and children out of props. The other properties,
such as id and data-topic, stay in rest.
\texttt{<section \{.\allowbreak{}.\allowbreak{}.\allowbreak{}rest\}>}
passes them to section. Only pass valid HTML properties to an HTML
element. Take out properties that are only meant for your own component.
TS uses \texttt{ComponentProps<'section'>} for the section's types.

If a spread and an explicit property use the same name, the last value
wins. RestFirst already takes title out of rest, so those two values
cannot conflict there. RestLast uses a separate heading property. This
lets us test what happens when rest still contains title: the later
spread replaces the earlier title.

\subsection{P5: named slots or
composition}\label{p5-named-slots-or-composition}

Slots takes header, footer and children. Header and footer are normal
named props. Children is the body between the tags. Slots creates the
header/main/footer layout. I would use this version when each card
should have the same structure.

Composition takes children and puts it inside article. The caller
chooses the elements and their order. I would use it when cards need
different structures. Both versions produce the same HTML in this
example. They give control of the layout to different parts of the code.

These are the functions. Both source files include the full JSX calls:

\begin{minipage}{\linewidth}

\begin{lstlisting}
function Slots({ header, footer, children }) {
  return <article><header>{header}</header><main>{children}</main><footer>{footer}</footer></article>
}
function Composition({ children }) {
  return <article>{children}</article>
}
\end{lstlisting}

\end{minipage}\par

\subsection{P6: undefined triggers a default, null does
not}\label{p6-undefined-triggers-a-default-null-does-not}

\texttt{function DefaultCard(\{ title = 'Free practice' \})} gives title
a default value. Undefined displays
\texttt{<h2>Free practice</\allowbreak{}h2>}. Null displays
\texttt{<h2></\allowbreak{}h2>}. A default is used for undefined or when
no value was given. It is not used for null. The TS type
\texttt{title?:\allowbreak{} string | null} allows both tests. To
replace both null and undefined, use \texttt{title ?? 'Free practice'}
instead.

\subsection{P7: three useful props
types}\label{p7-three-useful-props-types}

\texttt{\{ title:\allowbreak{} string; children?:\allowbreak{} ReactNode \}}
requires title. Children is optional. It accepts text, numbers,
elements, null and booleans. It does not accept a normal function child
or a plain object such as \texttt{\{ topic:\allowbreak{} 'read' \}}. A
function child needs its own function type, as in RenderCard.

\texttt{PropsWithChildren<\{ title:\allowbreak{} string \}>} adds
optional ReactNode children. Title is still required. Children does not
become required. Our Composition uses \texttt{PropsWithChildren} without
any other required props.

\texttt{ComponentProps<'button'>} gives us the usual button prop types.
It includes children, disabled, type and event handlers. TypedButton
takes children out and passes the other props to button. It accepts
type=``button'' and disabled. It rejects href because href belongs on a
link. The HTML output shows disabled. Testing a click handler would need
a browser.

\texttt{compiler-diagnostics.\allowbreak{}txt} has the actual rejected
cases: TS2322 for a function child, TS2353 for an object child, TS2322
for href on a button, and TS2710 for children written twice. These
errors are tested separately. We did not add more suppression comments
to the normal project.

\Needspace{12\baselineskip}

Types and functions from ts/props.tsx:

\begin{minipage}{\linewidth}

\begin{lstlisting}
type CardProps = { title: string; children?: ReactNode }
function Composition({ children }: PropsWithChildren) {
  return <article>{children}</article>
}
function TypedButton({ children, ...rest }: ComponentProps<'button'>) {
  return <button {...rest}>{children}</button>
}
\end{lstlisting}

\end{minipage}\par

\subsection{Actual React output
(JavaScript)}\label{actual-react-output-javascript}

\begin{minipage}{\linewidth}

\begin{lstlisting}
P1 object: <section><h2>Study</h2>Closures</section>
P1 pattern: <section><h2>Study</h2>Closures</section>
P2 nested: <section><h2>Study</h2>Practice</section>
P2 attribute: <section><h2>Study</h2>Practice</section>
P2 both: <section><h2>Study</h2>Nested</section>
P3 string: <section><h2>Study</h2>Read</section>
P3 string zero: <section><h2>Study</h2>0</section>
P3 numeric zero: <section><h2>Study</h2>0</section>
P3 element: <section><h2>Study</h2><strong>Code</strong></section>
P3 several: <section><h2>Study</h2><b>Read</b><i>Code</i></section>
P3 undefined: <section><h2>Study</h2></section>
P3 null: <section><h2>Study</h2></section>
P3 false: <section><h2>Study</h2></section>
P3 diagnostic: Functions are not valid as a React child. This may happen if you return %s instead of <%s /> from render. Or maybe you meant to call this function rather than return it. children children
P3 raw function: <section><h2>Study</h2></section>
P3 invoked function: <section><b>25 min</b></section>
P4 separated: <section id="s1" data-topic="props" title="Session">Read</section>
P4 explicit last: <section title="Explicit">Read</section>
P4 spread last: <section title="Spread">Read</section>
P5 slots: <article><header><h2>Plan</h2></header><main><p>Read</p></main><footer><small>Done</small></footer></article>
P5 composition: <article><header><h2>Plan</h2></header><main><p>Read</p></main><footer><small>Done</small></footer></article>
P6 undefined: <h2>Free practice</h2>
P6 null: <h2></h2>
P7 JS: No static props type checking; see ts/props.tsx and compiler probes.
\end{lstlisting}

\end{minipage}\par

\subsection{Actual P7 TypeScript
output}\label{actual-p7-typescript-output}

\begin{minipage}{\linewidth}

\begin{lstlisting}
P7 ReactNode: <section><h2>Typed</h2>7</section>
P7 PropsWithChildren: <article><em>Review</em></article>
P7 ComponentProps: <button type="button" disabled="">Start</button>
\end{lstlisting}

\end{minipage}\par

\subsection{Extra React experiment: key and a separate
ID}\label{extra-react-experiment-key-and-a-separate-id}

React uses key to identify an item. A component cannot read it as a
normal prop. If the component also needs the ID, pass it with a
different name, such as itemId. Both versions run in
\texttt{npm run all}. The assistant ran this extra experiment; it is not
a new prediction or mistake from me.

\Needspace{12\baselineskip}

JavaScript:

\begin{minipage}{\linewidth}

\begin{lstlisting}
function KeyCard(props) {
  console.log('key JS visible fields:', Object.keys(props).sort().join(','))
  const { key, itemId, children } = props
  console.log('key JS read:', String(key))
  return <section data-item-id={itemId}>{children}</section>
}
// Rendered with renderToStaticMarkup:
<KeyCard key="session-1" itemId="session-1">Read</KeyCard>
\end{lstlisting}

\end{minipage}\par

The JS source saves React's warning in stdout, then restores
console.error. The \texttt{\%s} below comes from that raw warning.
Object.keys lists the normal enumerable properties. Key does not appear
in that list.

The typed version uses itemId:

\begin{minipage}{\linewidth}

\begin{lstlisting}
type KeyCardProps = { itemId: string; children?: ReactNode }
function KeyCard(props: KeyCardProps) {
  console.log('key TS visible fields:', Object.keys(props).sort().join(','))
  const { itemId, children } = props
  return <section data-item-id={itemId}>{children}</section>
}
// key is allowed by JSX, without being a field of KeyCardProps:
<KeyCard key="session-1" itemId="session-1">Read</KeyCard>
\end{lstlisting}

\end{minipage}\par

\Needspace{7\baselineskip}

Actual output:

\begin{minipage}{\linewidth}

\begin{lstlisting}
key JS visible fields: children,itemId
key JS diagnostic: %s: `key` is not a prop. Trying to access it will result in `undefined` being returned. If you need to access the same value within the child component, you should pass it as a different prop. (https://react.dev/link/special-props) KeyCard
key JS read: undefined
key JS output: <section data-item-id="session-1">Read</section>
key TS visible fields: children,itemId
key TS output: <section data-item-id="session-1">Read</section>
\end{lstlisting}

\end{minipage}\par

The separate tests in
\texttt{tools/\allowbreak{}check-probes.\allowbreak{}cjs} check the
types too. JSX accepts key on KeyCard. But adding key to an object with
type \texttt{ComponentProps<typeof KeyCard>} gives TS2353. Reading key
from that type gives TS2339. The full errors are in
\texttt{compiler-diagnostics.\allowbreak{}txt}. The normal project
passes \texttt{npm run check}.

The fix is to read itemId in our code and let React use key. See
\href{https://react.dev/warnings/special-props}{React's special-props
explanation}.

\Needspace{27\baselineskip}

\section{JavaScript vs TypeScript}\label{javascript-vs-typescript}

JS and TS choose the same values when the code runs. These error
examples were tested separately. Their code and full messages are saved
in compiler-diagnostics.txt.

{\def\LTcaptype{none} % do not increment counter
\begin{longtable}[]{@{}
  >{\raggedright\arraybackslash}p{(\linewidth - 6\tabcolsep) * \real{0.2500}}
  >{\raggedright\arraybackslash}p{(\linewidth - 6\tabcolsep) * \real{0.2500}}
  >{\raggedright\arraybackslash}p{(\linewidth - 6\tabcolsep) * \real{0.2500}}
  >{\raggedright\arraybackslash}p{(\linewidth - 6\tabcolsep) * \real{0.2500}}@{}}
\toprule\noalign{}
\begin{minipage}[b]{\linewidth}\raggedright
Form/topic
\end{minipage} & \begin{minipage}[b]{\linewidth}\raggedright
Same when the code runs
\end{minipage} & \begin{minipage}[b]{\linewidth}\raggedright
TS type or check
\end{minipage} & \begin{minipage}[b]{\linewidth}\raggedright
Error in the wrong example
\end{minipage} \\
\midrule\noalign{}
\endhead
\bottomrule\noalign{}
\endlastfoot
1, tuple return & Read values by position & Return
\texttt{[number, string]} & TS2339: toUpperCase does not exist on
number \\
2, skipped array item & A missing item gives undefined & Check
\texttt{string \textbackslash{}| undefined} & TS18048: first is possibly
undefined \\
3, swap & Read the right side first & New values must fit the types &
TS2322: string is not assignable to number \\
6, rename & The colon gives a new name & Put the type after the pattern
& TS2339: toUpperCase does not exist on number \\
7, default & A default does not replace null & Allow null and check for
it & TS18047: minutes is possibly null \\
8, rest/source fields & Read properties the source has & Describe
available properties & TS2339: missing does not exist on the source
type \\
9, nesting & The inner object must exist & Check whether room exists &
TS2339: label does not exist on room-or-undefined \\
10, computed key & Use the property name in key & Use keyof for valid
property names & TS2322: unknown is not assignable to minutes \\
11, options default & A missing argument uses \{\} & Fields must allow
\texttt{\{\}} & TS2741: title is missing in type \{\} \\
12, entry loop & Each loop step gives a pair & Record and Map value
types & TS2339: toUpperCase does not exist on number \\
P7, children & Props is still one object & ReactNode or a
render-function signature & TS2322: function is not assignable to
ReactNode \\
\end{longtable}
}

\section{Verification and review}\label{verification-and-review}

The checks used Node.js 22.23.3 and npm 10.9.9. \texttt{npm run all} and
\texttt{npm run check} passed. Output.txt has the real run output.
Compiler-diagnostics.txt has expected errors from separate wrong-code
examples. Run
\texttt{node tools/\allowbreak{}check-probes.\allowbreak{}cjs} to repeat
those examples.

The log records my real guesses, surprises and later answers. It has
more than five wrong guesses or surprises. It also says which answers I
got right. It shows where I used AI help and where the assistant ran the
code. I still need to explain the guide in my own words. All four
bonuses are included. The public repository link is below.

\section{Bonus 1: LaTeX}\label{bonus-1-latex}

\texttt{latex/\allowbreak{}main.\allowbreak{}tex} and
\texttt{latex/\allowbreak{}main.\allowbreak{}pdf} contain the full
guide. They include the twelve forms, pitfalls, React examples,
comparison table, bonuses and personal log. The PDF is not a shorter
summary. It uses listings for code, real LaTeX tables, and a contents
page with links. In PDF code blocks, {[}U+1F9E0{]} stands for the brain
emoji. The source files and output files keep the original emoji.

Pandoc 3.14 converts this README to LaTeX. Tectonic 0.17.0 compiles it
to PDF. After installing these tools, run these commands in assignment4:

\begin{minipage}{\linewidth}

\begin{lstlisting}[language=sh]
pandoc README.md --standalone --shift-heading-level-by=-1 --toc --number-sections --listings --pdf-engine=xelatex --lua-filter=latex/layout.lua -V documentclass=article -V papersize=a4 -V fontsize=10pt -V geometry:margin=22mm -V title="Destructuring Field Guide" -V author="Tuba Nazl(*@\texttt{ı}@*) (*@\texttt{İ}@*)peky(*@\texttt{ü}@*)z - 2200004532" -V date="" -H latex/style.tex -o latex/main.tex
tectonic --keep-logs --keep-intermediates --outdir latex latex/main.tex
\end{lstlisting}

\end{minipage}\par

Pandoc only writes the TeX source here. Its pdf-engine option chooses a
preamble that supports Unicode. Tectonic then creates the PDF with its
XeTeX-based engine. The first Tectonic run may download the required TeX
packages. VERIFICATION.md records the build and visual checks. The LaTeX
editor can also compile main.tex. With both tools on PATH,
\texttt{npm run build:\allowbreak{}guide} runs these steps and adds the
name and student-number comment to the generated source.

\section{Bonus 4: Git repository}\label{bonus-4-git-repository}

\url{https://github.com/naz051223/assignment4-destructuring-field-guide}

Owner: \href{https://github.com/naz051223}{naz051223}. The repository is
public, as requested by the instructor. It can be read without a
collaborator invitation. The final upload was checked against the local
version.

The .gitignore excludes node\_modules. The real history shows more than
five meaningful commits across three days, including the final React
review. No earlier commit was redated or replaced. The latest runnable
version passed a clean installation, npm run all and npm run check. The
validation snapshot is recorded in CLEAN-CLONE-CHECK.txt.

\section{Bonus 2: Depth}\label{bonus-2-depth}

These extra examples test how destructuring works. The full code is in
\texttt{js/\allowbreak{}depth.\allowbreak{}mjs} and
\texttt{ts/\allowbreak{}depth.\allowbreak{}ts}. Run
\texttt{npm run depth}. Depth-output.txt has the real output.
\texttt{npm run all} includes the same examples. The assistant ran these
experiments. They are separate from my personal log.

\subsection{D1: an array pattern consumes an
iterator}\label{d1-an-array-pattern-consumes-an-iterator}

An array pattern reads values through an iterator: a way to get one
value at a time. Strings, Set, Map and generators all provide one. A Set
removes duplicates and keeps insertion order. A Map gives key-value
pairs. The nested pattern in our example takes the first pair, then its
key and value.

The string \texttt{'[U+1F9E0]A'} gives two values when iterated: the
emoji and A. Its length is 3 because string length counts UTF-16 code
units, and this emoji uses two units. Iteration counts code points here.
It does not always separate text into visible characters: a visible
character can contain more than one code point.

The generator logs start and closed. It does not reach the second yield.
The pattern asks for one value, then closes the iterator, which runs
finally. A rest pattern would keep reading the remaining values. With an
infinite generator, collecting all remaining values would not finish.

The plain object with index 0 and length 1 has no iterator. Its array
pattern throws a TypeError. Having numbered properties is not enough.

\subsection{D2: defaults are lazy and dependencies have an
order}\label{d2-defaults-are-lazy-and-dependencies-have-an-order}

The default function increases a counter. Values 0 and null keep their
values; undefined calls the function. The result is
\texttt{0 30 null 1}: one call, not three.

In \texttt{[work = 40, pause = work /\allowbreak{} 5]}, work is ready
before pause uses it. The result is 40 and 8. Reading a later variable
before it is initialized throws a ReferenceError.

The getter example records when each property is read. The pattern asks
for title first, so the order is title, minutes, even though the object
lists minutes first. The minutes getter returns undefined, so its
default runs. The final counter is 2 because the earlier experiment
already called the function once.

\subsection{D3: assignments, catch, and loop
inputs}\label{d3-assignments-catch-and-loop-inputs}

\texttt{[destination.\allowbreak{}minutes, destination.\allowbreak{}pause] = [40, 8]}
changes existing properties. It does not declare new variables. The
leading semicolon keeps the array assignment separate from the previous
statement.

\texttt{catch (\{ message \})} reads message from the Error object we
throw. It is not safe for every possible thrown value: JavaScript can
also throw null. For unknown errors, check the value before taking its
properties.

Our for\ldots in loop gives the array's index keys, 0 and 1. The
for\ldots of loop gives its values, arrays and props. Choose the loop
based on which values you need.

\subsection{D4: type-level extraction is not runtime
destructuring}\label{d4-type-level-extraction-is-not-runtime-destructuring}

First and Rest work on types. First finds the first type in a non-empty
tuple; Rest finds the remaining types. The keyword infer names the part
that was found. The readonly pattern also accepts tuples created with as
const.

For our tuple, First is the type 40 and Rest is
\texttt{['closures', true]}. The resulting Rest type is not readonly in
this example. These types check our variables; they do not create or cut
up an array at runtime.

\texttt{First<[]>} and \texttt{First<string[]>} are never here: neither
promises a first item. \texttt{Rest<[]>} is {[}{]}. This helper is for
tuples; ordinary arrays need a different rule.
\texttt{Parameters<typeof planSession>} gets the function's parameter
types:
\texttt{[topic:\allowbreak{} string, minutes:\allowbreak{} number]}. The
spread passes our two values to the function. TypeScript removes these
type rules before execution.

\subsection{D5: a readability refactoring with a defined behavior
change}\label{d5-a-readability-refactoring-with-a-defined-behavior-change}

The compact function reads several nested values in one pattern. It
works for the complete plan, but throws when owner is missing. The
longer version uses separate variables and checks missing objects first.
It returns Incomplete plan for missing data.

Both print Tuba: 40 min for the complete plan. They behave differently
for missing data, so the rewrite changes error handling as well as
readability. These small runs do not measure performance, so I do not
claim either is faster.

\subsection{Complete runnable JavaScript
investigation}\label{complete-runnable-javascript-investigation}

\begin{lstlisting}
// Tuba Nazl(*@\texttt{ı}@*) (*@\texttt{İ}@*)peky(*@\texttt{ü}@*)z (*@\texttt{—}@*) 2200004532
// Bonus 2: investigate iterator behavior, evaluation order and readability.
const [symbol, letter] = '[U+1F9E0]A'
console.log('depth string:', symbol, letter, '[U+1F9E0]A'.length)
const [firstUnique, ...otherUnique] = new Set(['arrays', 'arrays', 'props'])
console.log('depth Set:', firstUnique, JSON.stringify(otherUnique))
const [[day, minutes]] = new Map([['Monday', 40]])
console.log('depth Map:', day, minutes)
function* planned() {
  try {
    console.log('depth generator: start')
    yield 'arrays'
    console.log('depth generator: second yield reached')
    yield 'props'
  } finally {
    console.log('depth generator: closed')
  }
}
const [onlyFirst] = planned()
console.log('depth generator result:', onlyFirst)
// This is array-like but not iterable: an array pattern requires an iterator.
try {
  const [value] = { 0: 'arrays', length: 1 }
  console.log(value)
} catch (error) {
  console.log('depth array-like:', error.name + ': ' + error.message)
}

let calls = 0
function fallback() { calls++; return 30 }
const [zero = fallback(), missing = fallback(), explicitNull = fallback()] = [0, undefined, null]
console.log('depth lazy defaults:', zero, missing, explicitNull, calls)
const [work = 40, pause = work / 5] = []
console.log('depth earlier default:', work, pause)
try {
  const [early = later, later = 2] = []
  console.log(early, later)
} catch (error) {
  console.log('depth later default:', error.name + ': ' + error.message)
}
const reads = []
const settings = {
  get minutes() { reads.push('minutes'); return undefined },
  get title() { reads.push('title'); return 'React' }
}
const { title, minutes: duration = fallback() } = settings
console.log('depth getter order:', JSON.stringify(reads), title, duration, calls)

const destination = { minutes: 0, pause: 0 }
;[destination.minutes, destination.pause] = [40, 8]
console.log('depth targets:', JSON.stringify(destination))
try { throw new Error('Session unavailable') }
catch ({ message }) { console.log('depth catch:', message) }
const topics = ['arrays', 'props']
for (const index in topics) console.log('depth for-in:', index)
for (const topic of topics) console.log('depth for-of:', topic)

// A compact nested pattern can fail at an intermediate missing object.
function compact(plan) {
  const { owner: { name: student }, sessions: [{ minutes: duration }] } = plan
  return `${student}: ${duration} min`
}
function readable(plan) {
  const owner = plan.owner
  const firstSession = plan.sessions?.[0]
  if (owner == null || firstSession == null) return 'Incomplete plan'
  return `${owner.name}: ${firstSession.minutes} min`
}
const completePlan = { owner: { name: 'Tuba' }, sessions: [{ minutes: 40 }] }
console.log('depth compact:', compact(completePlan))
console.log('depth readable:', readable(completePlan))
try { console.log(compact({ sessions: [] })) }
catch (error) { console.log('depth compact missing:', error.name + ': ' + error.message) }
console.log('depth readable missing:', readable({ sessions: [] }))
\end{lstlisting}

\subsection{Complete TypeScript
investigation}\label{complete-typescript-investigation}

\begin{minipage}{\linewidth}

\begin{lstlisting}
// Tuba Nazl(*@\texttt{ı}@*) (*@\texttt{İ}@*)peky(*@\texttt{ü}@*)z (*@\texttt{—}@*) 2200004532
// Bonus 2: type-level tuple extraction, distinct from runtime destructuring.
export {}
type First<T extends readonly unknown[]> = T extends readonly [infer Head, ...unknown[]] ? Head : never
type Rest<T extends readonly unknown[]> = T extends readonly [unknown, ...infer Tail] ? Tail : []
const session = [40, 'closures', true] as const
const first: First<typeof session> = session[0]
const rest: Rest<typeof session> = [session[1], session[2]]
console.log('depth TS infer:', first, JSON.stringify(rest))
function planSession(topic: string, minutes: number): string {
  return `${topic}: ${minutes} min`
}
const argumentsForPlan: Parameters<typeof planSession> = ['props', 25]
console.log('depth TS Parameters:', planSession(...argumentsForPlan))
// An ordinary possibly-empty array does not guarantee a first tuple position.
type EmptyFirst = First<[]>
type ArrayFirst = First<string[]>
// Both aliases are never for this non-empty-tuple utility. No value is invented.
const emptyRest: Rest<[]> = []
console.log('depth TS empty Rest:', JSON.stringify(emptyRest))
// Keep the limitation explicit; this utility is for tuples, not arbitrary arrays.
\end{lstlisting}

\end{minipage}\par

\subsection{Actual depth output}\label{actual-depth-output}

\begin{minipage}{\linewidth}

\begin{lstlisting}
> destructuring-guide@0.0.0 depth
> node js/depth.mjs && tsx ts/depth.ts

depth string: [U+1F9E0] A 3
depth Set: arrays ["props"]
depth Map: Monday 40
depth generator: start
depth generator: closed
depth generator result: arrays
depth array-like: TypeError: {(intermediate value)(intermediate value)} is not iterable
depth lazy defaults: 0 30 null 1
depth earlier default: 40 8
depth later default: ReferenceError: Cannot access 'later' before initialization
depth getter order: ["title","minutes"] React 30 2
depth targets: {"minutes":40,"pause":8}
depth catch: Session unavailable
depth for-in: 0
depth for-in: 1
depth for-of: arrays
depth for-of: props
depth compact: Tuba: 40 min
depth readable: Tuba: 40 min
depth compact missing: TypeError: Cannot read properties of undefined (reading 'name')
depth readable missing: Incomplete plan
depth TS infer: 40 ["closures",true]
depth TS Parameters: props: 25 min
depth TS empty Rest: []
\end{lstlisting}

\end{minipage}\par

References:
\href{https://tc39.es/ecma262/multipage/ecmascript-language-statements-and-declarations.html\#sec-destructuring-binding-patterns-runtime-semantics-iteratorbindinginitialization}{ECMAScript
iterator binding initialization},
\href{https://www.typescriptlang.org/docs/handbook/2/conditional-types.html}{TypeScript
conditional types},
\href{https://www.typescriptlang.org/docs/handbook/utility-types.html\#parameterstype}{Parameters
utility}.

\section{Bonus 3: Real code}\label{bonus-3-real-code}

These five cases use three sources: the instructor's Assignment 3
starter (\texttt{Downloads/\allowbreak{}starter.\allowbreak{}zip}),
TypeScript 5.8.3, and TanStack Query. The TanStack version is commit
\texttt{36dd7927c85ce6cee37ab1eb3af25e7fc7ccbb14}. The line numbers
refer to those exact versions. References/source-excerpts.txt keeps the
source extracts and SHA-256 hashes that identify the files. The full
outside projects are not included. Each rewrite shows the same selected
code without destructuring. These are code comparisons, not separate
apps. Hooks stay at the component's top level. Do not call the hook
again just to read a second value.

\subsection{Case 1: TaskRow receives one props
object}\label{case-1-taskrow-receives-one-props-object}

Source: Assignment 3 starter,
\texttt{starter/\allowbreak{}src/\allowbreak{}TaskRow.\allowbreak{}tsx},
line 5.

\begin{minipage}{\linewidth}

\begin{lstlisting}
export default function TaskRow({ task, onNext, onDelete }: TaskRowProps) {
  return (
    <li style={{ display: 'flex', gap: 8, alignItems: 'center' }}>
      <span style={{ flex: 1 }}>{task.title}</span>
      <em>{task.status}</em>
      <button onClick={onNext}>Next</button>
      <button onClick={onDelete}>Delete</button>
    </li>
  )
}
\end{lstlisting}

\end{minipage}\par

Forms: 5 and 11. Expanded parameter access:

\begin{minipage}{\linewidth}

\begin{lstlisting}
export default function TaskRow(props: TaskRowProps) {
  const task = props.task
  const onNext = props.onNext
  const onDelete = props.onDelete
  return (
    <li style={{ display: 'flex', gap: 8, alignItems: 'center' }}>
      <span style={{ flex: 1 }}>{task.title}</span>
      <em>{task.status}</em>
      <button onClick={onNext}>Next</button>
      <button onClick={onDelete}>Delete</button>
    </li>
  )
}
\end{lstlisting}

\end{minipage}\par

TaskList passes task, onNext and onDelete as named props. TaskRowProps
has no children field, and TaskRow does not read children. Next and
Delete are children of the inner buttons. They are not children passed
to TaskRow. React uses key=\{t.id\} to identify list items; key is not
an ordinary props field. I would keep this short pattern because it
shows which props the component uses.

\subsection{Case 2: useReducer returns a
pair}\label{case-2-usereducer-returns-a-pair}

Source: Assignment 3 starter,
\texttt{starter/\allowbreak{}src/\allowbreak{}App.\allowbreak{}tsx},
line 37.

\begin{minipage}{\linewidth}

\begin{lstlisting}
const [tasks, dispatch] = useReducer(taskReducer, [])
\end{lstlisting}

\end{minipage}\par

Form: 1. Expanded:

\begin{minipage}{\linewidth}

\begin{lstlisting}
const reducerResult = useReducer(taskReducer, [])
const tasks = reducerResult[0]
const dispatch = reducerResult[1]
\end{lstlisting}

\end{minipage}\par

Tasks holds the current state. Dispatch is the function used to send
actions. The first and second positions decide what we receive, not the
variable names. The rewrite calls useReducer once and reads both results
from that call. Calling it twice would create two separate hooks. Named
imports also use braces, but their syntax is different from object
destructuring.

\subsection{Case 3: a library helper extracts pattern
fields}\label{case-3-a-library-helper-extracts-pattern-fields}

Source:
\texttt{node\_modules/\allowbreak{}typescript/\allowbreak{}lib/\allowbreak{}typescript.\allowbreak{}js},
version 5.8.3, lines 3631--3633; same function is visible in the
\href{https://unpkg.com/typescript@5.8.3/lib/typescript.js}{versioned
package}.

\begin{minipage}{\linewidth}

\begin{lstlisting}
function patternText({ prefix, suffix }) {
  return `${prefix}*${suffix}`;
}
\end{lstlisting}

\end{minipage}\par

Forms: 5 and 11. Expanded:

\begin{minipage}{\linewidth}

\begin{lstlisting}
function patternText(pattern) {
  const prefix = pattern.prefix
  const suffix = pattern.suffix
  return `${prefix}*${suffix}`
}
\end{lstlisting}

\end{minipage}\par

The function puts * between prefix and suffix. Both values come from one
object, not two separate arguments. Both versions fail if that object is
null or undefined. I would keep the original pattern because two fields
are easy to read.

\subsection{Case 4: a dense library pattern I would
simplify}\label{case-4-a-dense-library-pattern-i-would-simplify}

Source: TypeScript 5.8.3,
\texttt{node\_modules/\allowbreak{}typescript/\allowbreak{}lib/\allowbreak{}typescript.\allowbreak{}js},
line 4872, inside tryParseComponents. The preceding lines check that the
regular expression match exists.

\begin{minipage}{\linewidth}

\begin{lstlisting}
const [, major, minor = "0", patch = "0", prerelease = "", build2 = ""] = match;
\end{lstlisting}

\end{minipage}\par

Forms: 2 (skip), 4 (defaults). Expanded:

\begin{minipage}{\linewidth}

\begin{lstlisting}
const major = match[1]
const minor = match[2] === undefined ? '0' : match[2]
const patch = match[3] === undefined ? '0' : match[3]
const prerelease = match[4] === undefined ? '' : match[4]
const build2 = match[5] === undefined ? '' : match[5]
\end{lstlisting}

\end{minipage}\par

The first match item is the whole matched text. The leading comma skips
it. Missing groups use defaults. I find this pattern harder to read
because I have to count positions and check the regular expression to
understand each value. In code I maintain, I would use named groups or a
small helper with clear names. BuildMetadata would also be clearer than
build2. This is generated library code, so I would not edit the
installed file itself. The rewrite uses the normal match array returned
by this function's regular expression.

\subsection{Case 5: an open-source hook
result}\label{case-5-an-open-source-hook-result}

Source: TanStack Query,
\href{https://github.com/TanStack/query/blob/36dd7927c85ce6cee37ab1eb3af25e7fc7ccbb14/examples/react/simple/src/index.tsx\#L21-L30}{examples/react/simple/src/index.tsx,
lines 21--30}.

\begin{minipage}{\linewidth}

\begin{lstlisting}
const { isPending, error, data, isFetching } = useQuery({
  queryKey: ['repoData'],
  queryFn: async () => {
    const response = await fetch(
      'https://api.github.com/repos/TanStack/query',
    )
    if (!response.ok) throw new Error('Failed to fetch repo data')
    return await response.json()
  },
})
\end{lstlisting}

\end{minipage}\par

Form: 5. Expanded access, preserving one hook call:

\begin{minipage}{\linewidth}

\begin{lstlisting}
const query = useQuery({
  queryKey: ['repoData'],
  queryFn: async () => {
    const response = await fetch(
      'https://api.github.com/repos/TanStack/query',
    )
    if (!response.ok) throw new Error('Failed to fetch repo data')
    return await response.json()
  },
})
const isPending = query.isPending
const error = query.error
const data = query.data
const isFetching = query.isFetching
\end{lstlisting}

\end{minipage}\par

These fields come from the object returned by useQuery, not from props.
The rewrite keeps the same options and calls the hook once. The
component checks loading and errors before showing data, and uses
isFetching for Updating\ldots. The same values are selected at runtime.
In typed code, separate variables can lose the relationship TypeScript
uses to check a result after a condition. Keeping query.isPending and
query.data on one object can preserve that relationship. I would keep
the original short pattern. This is a source-code comparison; I did not
run its GitHub network request as part of the local examples.

\section{Extra example: const and
readonly}\label{extra-example-const-and-readonly}

Run \texttt{npm run readonly} to check these assignments with
TypeScript. The command is also part of \texttt{npm run all}. It checks
each line separately, so one expected error does not hide the results of
the other lines.

\begin{minipage}{\linewidth}

\begin{lstlisting}
type Session = {
  readonly minutes: number;
  readonly room: { label: string };
};
const session: Session = { minutes: 40, room: { label: 'A' } };
const { minutes, room } = session;
\end{lstlisting}

\end{minipage}\par

\texttt{const} prevents assigning a new value to the local variable
\texttt{minutes}. \texttt{readonly} prevents assigning a new value to
\texttt{session.\allowbreak{}minutes} or
\texttt{session.\allowbreak{}room}. However, \texttt{label} is not
readonly, so \texttt{session.\allowbreak{}room.\allowbreak{}label} can
still change.

Destructuring \texttt{room} gives us another reference to the same
object. Changing its label changes the shared object. A number works
differently: a new \texttt{let} variable can change without changing the
number stored in \texttt{session}.

Actual output, also saved in \texttt{readonly-output.\allowbreak{}txt}:

\begin{minipage}{\linewidth}

\begin{lstlisting}
minutes = 60;
TS2588: Cannot assign to 'minutes' because it is a constant.
room.label = 'B';
OK (no compiler errors)
session.minutes = 60;
TS2540: Cannot assign to 'minutes' because it is a read-only property.
session.room = { label: 'B' };
TS2540: Cannot assign to 'room' because it is a read-only property.
session.room.label = 'C';
OK (no compiler errors)
let { minutes: editableMinutes } = session; editableMinutes = 60;
OK (no compiler errors)
runtime: 40 60 40 C
runtime frozen: false false
\end{lstlisting}

\end{minipage}\par

The runtime values are 40 for the const number, 60 for the local let
number, 40 for the original number, and C for the shared label. Only
accepted code runs. \texttt{readonly} checks assignments in TypeScript;
it does not freeze objects at runtime. Both frozen checks are false. To
protect label too, use \texttt{readonly label:\allowbreak{} string}.

\section{Extra example: changing an object or replacing a
reference}\label{extra-example-changing-an-object-or-replacing-a-reference}

Run \texttt{npx tsx ts/\allowbreak{}reference-followup.\allowbreak{}ts}.
The experiment is also included in \texttt{npm run ts} and
\texttt{npm run all}, and is checked by \texttt{npm run check}. The
following is the complete runnable source:

\begin{minipage}{\linewidth}

\begin{lstlisting}
// Tuba Nazl(*@\texttt{ı}@*) (*@\texttt{İ}@*)peky(*@\texttt{ü}@*)z (*@\texttt{—}@*) 2200004532
// Changing a shared object versus assigning a new object.
export {}

{
  const session: { readonly room: { label: string } } = {
    room: { label: 'A' }
  };
  let { room } = session;
  console.log('shared before:', room === session.room);
  room.label = 'B';
  room = { label: 'C' };
  console.log('shared after:', room === session.room);
  console.log('after rebinding:', room.label, session.room.label);
  room.label = 'D';
  console.log('after local mutation:', room.label, session.room.label);
}

{
  const original = { label: 'A' };
  let copy = original;
  copy = { label: 'B' };
  original.label = 'C';
  console.log('ordinary alias:', copy.label, original.label);
}
\end{lstlisting}

\end{minipage}\par

Actual output, also saved in
\texttt{reference-followup-output.\allowbreak{}txt}:

\begin{minipage}{\linewidth}

\begin{lstlisting}
shared before: true
shared after: false
after rebinding: C B
after local mutation: D B
ordinary alias: B C
\end{lstlisting}

\end{minipage}\par

At first, \texttt{room} and \texttt{session.\allowbreak{}room} point to
the same object. This is why \texttt{room === session.\allowbreak{}room}
is true. Setting \texttt{room.\allowbreak{}label} to B changes that
object.

Then \texttt{room = \{ label:\allowbreak{} 'C' \}} makes the local
variable point to a new object. \texttt{session.\allowbreak{}room} still
points to the first object, whose label is B. The two references are now
different. Changing the new object's label to D leaves the first object
at B. This gives C B first, then D B.

The second example uses the same rule without destructuring. Assigning a
new object to \texttt{copy} leaves it separate from \texttt{original}.
The final labels are B and C. The name \texttt{copy} does not create a
separate object by itself.

\section{Sources and assistance}\label{sources-and-assistance}

The starter and handout gave the project structure. I used AI help for
the examples and code. The log separates my guesses from the assistant's
actual runs on this computer. These sources were used to check the
rules, not to copy examples:

\begin{itemize}
\tightlist
\item
  \href{https://developer.mozilla.org/en-US/docs/Web/JavaScript/Reference/Operators/Destructuring}{MDN:
  Destructuring}: defaults and source-value rules.
\item
  \href{https://www.typescriptlang.org/docs/handbook/2/objects.html}{TypeScript
  Handbook: Object Types}: property patterns, optional fields, tuples
  and readonly tuples.
\item
  \href{https://react.dev/learn/passing-props-to-a-component}{React:
  Passing Props to a Component}: one props object, children and
  defaults.
\item
  \href{https://react.dev/reference/react-dom/components/common}{React
  DOM: Common components}: renderable children values.
\item
  \href{https://legacy.reactjs.org/docs/render-props.html}{React: Render
  Props}: explicitly called function props.
\end{itemize}

\section{Submission contents}\label{submission-contents}

The ZIP contains one folder named assignment4. It does not include
node\_modules. It includes the guide, code, package files,
tsconfig.json, output.txt, compiler-diagnostics.txt, LaTeX source and
PDF. After extracting it, run npm install to install the dependencies.
All four bonuses are included. The public repository link is above. The
GitHub upload was checked against the local version.

\section{My own log}\label{my-own-log}

These are my real guesses and corrections from the review with an AI
assistant. The assistant ran the code on this computer. I discussed the
examples and results. The notes include wrong guesses, surprises and
correct later answers. They do not claim that I found or ran every error
on my own. The exact error messages and outputs are above and in the
output files.

\subsection{1. I expected the default function to
run}\label{i-expected-the-default-function-to-run}

In \texttt{js/\allowbreak{}default-review.\allowbreak{}mjs}, I thought
fallback() would run. It did not, which surprised me. The assistant ran
the code on this computer. The output was \texttt{0 0}: minutes stayed 0
and the function was not called.

I did not know the default rule clearly before this example. A default
is used for undefined or when no value was given. It is not used for 0.
The default function only runs when the default is needed. There was no
error message.

\begin{minipage}{\linewidth}

\begin{lstlisting}
let calls = 0
function fallback() {
  calls++
  return 30
}
const { minutes = fallback() } = { minutes: 0 }
console.log(minutes, calls)
\end{lstlisting}

\end{minipage}\par

\Needspace{7\baselineskip}

Run \texttt{node js/\allowbreak{}default-review.\allowbreak{}mjs}.
Actual output:

\begin{minipage}{\linewidth}

\begin{lstlisting}
0 0
\end{lstlisting}

\end{minipage}\par

\subsection{2. I expected renaming to change the existing title
variable}\label{i-expected-renaming-to-change-the-existing-title-variable}

In \texttt{js/\allowbreak{}rename-review.\allowbreak{}mjs}, I guessed
that both lines would print React çalışması. I thought making the values
equal would also change title. The assistant ran it on this computer.
The actual output was:

\begin{minipage}{\linewidth}

\begin{lstlisting}
Haftal(*@\texttt{ı}@*)k plan
React (*@\texttt{ç}@*)al(*@\texttt{ı}@*)(*@\texttt{ş}@*)mas(*@\texttt{ı}@*)
\end{lstlisting}

\end{minipage}\par

\begin{minipage}{\linewidth}

\begin{lstlisting}
const title = 'Haftal(*@\texttt{ı}@*)k plan'
const session = { title: 'React (*@\texttt{ç}@*)al(*@\texttt{ı}@*)(*@\texttt{ş}@*)mas(*@\texttt{ı}@*)' }
const { title: heading } = session
console.log(title)
console.log(heading)
\end{lstlisting}

\end{minipage}\par

The title inside the pattern is the property name. Heading is the new
variable. The line means
\texttt{const heading = session.\allowbreak{}title}. It does not change
the existing title variable, so title stays Haftalık plan. There was no
error message. I needed to separate reading a property from changing
another variable. Run
\texttt{node js/\allowbreak{}rename-review.\allowbreak{}mjs}.

\subsection{3. I compared a number with a shared nested
object}\label{i-compared-a-number-with-a-shared-nested-object}

In \texttt{js/\allowbreak{}reference-review.\allowbreak{}mjs}, I
correctly guessed 40 and B. Minutes received 40 when destructuring ran.
The later change to session.minutes did not change that separate number
variable. Room and session.room still pointed to the same object, so
changing its label was visible through both names.

\begin{minipage}{\linewidth}

\begin{lstlisting}
const session = {
  minutes: 40,
  room: { label: 'A' }
}
const { minutes, room } = session
session.minutes = 60
room.label = 'B'
console.log(minutes)
console.log(session.room.label)
\end{lstlisting}

\end{minipage}\par

\Needspace{7\baselineskip}

The assistant ran this code on this computer. Actual output:

\begin{minipage}{\linewidth}

\begin{lstlisting}
40
B
\end{lstlisting}

\end{minipage}\par

This example helped me compare a number with an object reference.
Destructuring does not make a separate copy of every nested object.
There was no error message. Run
\texttt{node js/\allowbreak{}reference-review.\allowbreak{}mjs}.

\subsection{4. I overlooked the zero passed to a conditional
child}\label{i-overlooked-the-zero-passed-to-a-conditional-child}

I guessed that Çalışma var would appear. I then explained that I had
missed count=\{0\} in the example. I knew that \&\& stops at a falsy
left value, but I had not noticed which value was passed in.

\begin{minipage}{\linewidth}

\begin{lstlisting}
function Session({ count }) {
  return <section>{count && <span>(*@\texttt{Ç}@*)al(*@\texttt{ı}@*)(*@\texttt{ş}@*)ma var</span>}</section>
}
// Rendered usage: <Session count={0} />
\end{lstlisting}

\end{minipage}\par

The assistant rendered it on this computer. The actual HTML was
\texttt{<section>0</\allowbreak{}section>}. Here \&\& returns 0, not
false. React shows the number 0 as text.

\Needspace{7\baselineskip}

In \texttt{js/\allowbreak{}children-review.\allowbreak{}jsx}, the fixed
check is \texttt{count > 0 \&\& <span>Çalışma var</\allowbreak{}span>}.
For count=0, it returns false. React does not show false. There was no
error message. Run
\texttt{npx tsx js/\allowbreak{}children-review.\allowbreak{}jsx}.
Actual output:

\begin{minipage}{\linewidth}

\begin{lstlisting}
original: <section>0</section>
checked: <section></section>
\end{lstlisting}

\end{minipage}\par

\subsection{5. I traced a function child and its returned
element}\label{i-traced-a-function-child-and-its-returned-element}

I said children was a function and separated it from children(). That
was correct for this example. Children holds the function; children(2)
calls it. The call returns a React element. It is not another function.
I also correctly calculated 5 for number + 3. This note records an
example I checked, not a wrong answer.

\begin{minipage}{\linewidth}

\begin{lstlisting}
function Card({ children }) {
  return <section>{children(2)}</section>
}
// Rendered usage:
<Card>{number => <strong>{number + 3}</strong>}</Card>
\end{lstlisting}

\end{minipage}\par

\Needspace{7\baselineskip}

The assistant rendered it on this computer. Actual HTML:

\begin{minipage}{\linewidth}

\begin{lstlisting}[language=HTML]
<section><strong>5</strong></section>
\end{lstlisting}

\end{minipage}\par

React does not call the function automatically. Card calls it and
displays the result. The TS type must describe the function:
\texttt{\{ children:\allowbreak{} (number:\allowbreak{} number) => ReactNode \}}.
Ordinary ReactNode children do not include this function type. There was
no error message. Run
\texttt{npx tsx js/\allowbreak{}function-review.\allowbreak{}jsx}.

\subsection{6. I predicted the wrong length for a destructured string
property}\label{i-predicted-the-wrong-length-for-a-destructured-string-property}

For
\texttt{const \{ length \} = 'React'; console.\allowbreak{}log(length)},
I guessed 2. I explained that I had thought of React as one whole word.
I had not paid attention to its characters. The assistant ran it on this
computer. The output was 5.

The pattern reads the string's length property. It does not count words
or the variables in the pattern. JavaScript temporarily wraps the string
to read the property. React has five ASCII characters. In general,
length counts UTF-16 code units, so it may differ from the number of
visible characters. There was no error message. The code is in
\texttt{js/\allowbreak{}edge-review.\allowbreak{}mjs}.

\subsection{7. I expected a later default binding to be
available}\label{i-expected-a-later-default-binding-to-be-available}

For \texttt{const [a = b, b = 2] = []}, I guessed b for a and 2 for b.
The assistant ran the code on this computer. It did not print the
values. It threw this exact error:

\begin{minipage}{\linewidth}

\begin{lstlisting}
ReferenceError: Cannot access 'b' before initialization
\end{lstlisting}

\end{minipage}\par

The first missing item uses a's default straight away. At that point b
is not ready yet, so reading it throws. The code stops before b gets its
value 2. This is called the temporal dead zone: a let or const variable
cannot be read before it is initialized. Putting the earlier value first
fixes the dependency:

\begin{minipage}{\linewidth}

\begin{lstlisting}
const [b = 2, a = b] = []
console.log(a, b) // actual output: 2 2
\end{lstlisting}

\end{minipage}\par

\texttt{js/\allowbreak{}edge-review.\allowbreak{}mjs} includes the
caught error and the working example.

\subsection{8. I forgot what readonly
prevents}\label{i-forgot-what-readonly-prevents}

I correctly guessed that the const variable minutes could not be
assigned 60, while room.label could change. After readonly was added, I
thought all three assignments would work. I was surprised that assigning
to session.minutes and replacing session.room both gave TS2540 errors. I
had forgotten that readonly prevents assigning new values to those
properties.

The assistant checked the code on this computer.
\texttt{session.\allowbreak{}room.\allowbreak{}label = 'C'} was allowed
because label was not readonly. The extra example above has the actual
compiler messages and output.

\subsection{9. I mixed up changing an object and replacing a
reference}\label{i-mixed-up-changing-an-object-and-replacing-a-reference}

I guessed A and B for the room example. The result was C and B. After
\texttt{room.\allowbreak{}label = 'D'}, I guessed D and C. The result
was D and B. For the copy example, I guessed A and C. The result was B
and C. The assistant ran the examples on this computer and saved their
output.

I was mixing up two steps. \texttt{room.\allowbreak{}label = 'B'}
changes the shared object. \texttt{room = \{ label:\allowbreak{} 'C' \}}
makes only room point to a new object. The original object keeps its B
label.

My earlier answer of 40 and 60 in the number example was an order slip.
I said I had not paid attention to the order. It is not counted as a new
mistake about how the code works.

In a later example, I correctly guessed C and D. Room pointed to a new
object, so changing the original object's label did not change it. I
also explained the default rule correctly: a default is used for
undefined or when no value was given, but not for null. I corrected a
typing slip in that sentence.

\end{document}
